\documentclass[11pt,a4paper]{article}
\usepackage[margin=27mm]{geometry}
\usepackage{longtable,booktabs,array,fvextra}
\usepackage{amsmath,amsthm,amssymb,fontspec,unicode-math,hyperref}
\setmainfont{STIXGeneral}
\setmathfont{STIX Two Math}
\setmonofont{Menlo}
\hypersetup{hidelinks}
\usepackage{newunicodechar}
\newfontfamily\numeralfallback{STIX Two Text}
\newunicodechar{∎}{\ensuremath{\blacksquare}}
\newunicodechar{⁰}{{\numeralfallback\char"2070\relax}}
\newunicodechar{₀}{{\numeralfallback\char"2080\relax}}
\newunicodechar{₁}{{\numeralfallback\char"2081\relax}}
\newunicodechar{₂}{{\numeralfallback\char"2082\relax}}
\newunicodechar{₃}{{\numeralfallback\char"2083\relax}}
\newunicodechar{₄}{{\numeralfallback\char"2084\relax}}
\newunicodechar{⁻}{{\numeralfallback\char"207B\relax}}
\newunicodechar{⁸}{{\numeralfallback\char"2078\relax}}
\newunicodechar{ⁱ}{{\numeralfallback\char"2071\relax}}
\newunicodechar{⁴}{{\numeralfallback\char"2074\relax}}
\newunicodechar{⁷}{{\numeralfallback\char"2077\relax}}
\newunicodechar{⁶}{{\numeralfallback\char"2076\relax}}
\newunicodechar{⁵}{{\numeralfallback\char"2075\relax}}
\newunicodechar{₈}{{\numeralfallback\char"2088\relax}}
\newunicodechar{ᵢ}{\textsubscript{i}}
\newunicodechar{ⱼ}{\textsubscript{j}}
\newunicodechar{ₜ}{\textsubscript{t}}
\providecommand{\tightlist}{\setlength{\itemsep}{0pt}\setlength{\parskip}{0pt}}
\newtheorem{theorem}{Theorem}
\newtheorem{proposition}[theorem]{Proposition}
\newtheorem{lemma}[theorem]{Lemma}
\newtheorem{corollary}[theorem]{Corollary}
\theoremstyle{remark}\newtheorem{remark}[theorem]{Remark}
\newcommand{\gen}[1]{\langle #1\rangle}
\newcommand{\normalgen}[1]{\langle\!\langle #1\rangle\!\rangle}
\newcommand{\Z}{\mathbb Z}
\newcommand{\F}{\mathbb F}
\setlength{\emergencystretch}{2em}

\providecommand{\passthrough}[1]{#1}

\title{Supported kernels, translation lattices and finite marked quotients}
\author{Human Compression Project}
\date{}
\begin{document}
\maketitle
\begin{abstract}
We study marked groups through the intersections of subgroups generated
by subsets of their markers. A construction over a group with deletion
maps attaches local extensions whose kernels are invisible on proper
local subgroups. We prove generation of the full compatible product,
compute every parabolic kernel and obtain the intersection property. The
same coordinates give local recognition, deletion maps, lower-central
identities, separator factorizations and transfer of the flag property
for coset geometries. For finite two-groups whose noncommuting diagram
is a forest, we identify a canonical class-two quotient and a finite,
exhaustive parameter space of marked quotients of any prescribed excess
order. Explicit affine square groups supply local models of every
required two-power order. The paper includes the supporting quotient,
Frattini and amalgam arguments so that subsequent applications can cite
complete statements with specified hypotheses.
\end{abstract}
\section{Introduction}
A presentation by involutions does not by itself determine a regular
polytope or hypertope. The distinguished generators must have the
prescribed pair orders, their generated subgroups must intersect
correctly, and a geometric conclusion may require a further
common-representative condition on cosets. Establishing these properties
separately for every presentation can obscure the part of the proof that
is independent of its numerical parameters.

The main construction here retains an already understood marked quotient
and records the extra relations in separate kernel coordinates. Each
attached group is required to agree with the quotient on every proper
part of its support. This hypothesis lets us isolate its whole kernel,
even when one support contains another. Once the kernels have been
isolated, all marked subgroup intersections and orders follow from a
single coordinate calculation. The geometric assertion uses an
additional proved transfer of typed coset Helly; it is not inferred from
the intersection property alone.

A second branch identifies the quotient which every finite two-group
marking on a forest must have. The proof uses the intersection property
in the original group, a descending lower-central argument and the
separation of a forest at an edge. It then reduces classification at any
fixed excess order to normal subgroups of one finite nilpotent quotient.
The explicit square groups in the affine-square section provide a complementary
construction: their orders vary, while their proper marked subgroups
remain unchanged.

The companion applications use these mechanisms for prescribed two-power
orders, three-generator hypertopes, square-section assemblies, common
quotients, high-type classifications and the nondegenerate order-1024 census.
Construction and presentation recognition are distinct: the former isolates
local kernels in a compatible product, while the latter proves that a proposed
presentation has no further kernel. Both steps are proved below. The source
results retain their original attribution; no historical priority for these
shared arguments is asserted.
The same kernel language controls the mixed-prime affine families. An exact
Schreier presentation first identifies the infinite translation kernel; lattice
periods then determine the quotient and every marked intersection. We give a
criterion for arbitrary parabolics, including those with nonzero translation
kernels, before specializing to the faithful finite adjacent subgroups of the
rank-three applications. The resulting period-and-local-extension theorem
combines this recognition with the supported construction. The complete
square normal form and the kernel-isolation proof occur only once here.

This volume retains all seven previously admitted Topic A original contracts
and the complete mixed-prime contract L012. It adds no original or conjecture
credit. It preserves all source corrections, all eleven whole required
program/data files, and the separate retained classical and software inputs.
The arbitrary-parabolic lattice criterion is newly stated in this edition;
its elementary coset proof and its current lack of full Lean formalization are
explicit. All numerical comparisons remain deferred to a separately reviewed
whole-volume boundary.

\section{Markings, images and kernels}\label{markings-images-and-kernels}

A marked group is a group \(G\) together with a finite indexed
generating family \((r_i)_{i\in I}\). Except where stated, the markers
need not be involutions or distinct. For \(S\subseteq I\), put \[
G_S=\langle r_i:i\in S\rangle,\qquad G_\varnothing=1.
\] We call these subgroups parabolics, retaining their index sets even
when two of them happen to coincide. The marking has the intersection
property, or IP, if \[
G_S\cap G_T=G_{S\cap T}\quad(S,T\subseteq I).
\tag{M1}
\] A marked homomorphism sends each source marker to the corresponding
target marker. Its image on \(G_S\) is therefore the target parabolic
with index set \(S\).

All subgroup products below are products of sets unless a subgroup
assertion is made. We write \(\gamma_1(G)=G\) and
\(\gamma_{j+1}(G)=[\gamma_j(G),G]\). We use the elementary
correspondence and isomorphism theorems for groups. In the finite
\(p\)-group arguments we retain Burnside's basis theorem, the nilpotence
of finite \(p\)-groups, and the usual multilinear commutator description
of the lower-central graded Lie ring. Precisely, its degree-one part
generates the graded Lie ring, and degree \(j\) is generated by
left-normed commutators of weight \(j\) in any generating set. In
the amalgam section we explicitly retain the normal-form theorem for an
amalgamated free product. None of these classical inputs is presented as
a result originating in this manuscript.

\textbf{Lemma UA.K0 (lifting one intersection).} Let \(q:H\to Q\) be a
homomorphism and let \(C\leq A\cap B\), where \(A,B,C\leq H\). If \[
q(A)\cap q(B)=q(C),
\] then \[
A\cap B=C\quad\Longleftrightarrow\quad A\cap B\cap\ker q\leq C.
\tag{M2}
\]

\emph{Proof.} The forward implication is immediate. For the converse,
let \(x\in A\cap B\). Its image belongs to \(q(C)\), so choose
\(c\in C\) with \(q(c)=q(x)\). Then \(xc^{-1}\in A\cap B\cap\ker q\),
which is contained in \(C\) by hypothesis. Multiplying by \(c\) gives
\(x\in C\). The inclusion \(C\leq A\cap B\) gives equality. \(\square\)

\textbf{Theorem UA.K (kernel profiles).} Let \(q:H\twoheadrightarrow Q\)
be a marked epimorphism and suppose that the marking of \(Q\) has IP.
Write \(K=\ker q\) and \(K_S=K\cap H_S\). Then \(H\) has IP if and only
if \[
K_S\cap K_T=K_{S\cap T}\quad(S,T\subseteq I).
\tag{M3}
\]

\emph{Proof.} Apply UA.K0 to \(A=H_S\), \(B=H_T\), and
\(C=H_{S\cap T}\). The image hypothesis is exactly IP in \(Q\). The
subgroup on the left of the containment in (M2) is \(K_S\cap K_T\). Its
containment in \(H_{S\cap T}\) is equivalent to its containment in
\(K_{S\cap T}\); the reverse inclusion always holds. This proves each
required intersection and its converse. \(\square\)

The next result concerns the opposite direction. Passing to a quotient
can create new intersections, and this defect is a product condition in
its kernel.

\textbf{Theorem UA.Q (exact intersections of images).} Let \(q:H\to Q\)
be a homomorphism with kernel \(N\), and let \(A,B\leq H\). Then \[
q(A)\cap q(B)=q(A\cap B)
\quad\Longleftrightarrow\quad
N\cap AB=(N\cap A)(N\cap B).
\tag{M4}
\] If \(q\) is injective on \(A\) and \(B\), the right side reduces to
\(N\cap AB=\{1\}\). If \(N\leq A\), the equality always holds.

\emph{Proof.} Suppose the image intersection is exact, and take
\(ab\in N\), where \(a\in A\) and \(b\in B\). Since \(q(a)=q(b)^{-1}\),
choose \(c\in A\cap B\) with this image. Then \[
ab=(ac^{-1})(cb)\in(N\cap A)(N\cap B).
\] The reverse inclusion requires no hypothesis. Conversely, if
\(q(a)=q(b)\), then \(ab^{-1}\in N\cap AB\). Write \(ab^{-1}=uv\), with
\(u\in N\cap A\) and \(v\in N\cap B\). The element \(u^{-1}a=vb\) lies
in \(A\cap B\) and has image \(q(a)\), proving the required image
equality. Under injectivity both factors on the right of (M4) are
trivial. If \(N\leq A\), both sides of that equality are \(N\).
\(\square\)

\subsection*{The rank-three string test}
\textbf{Rank-three test.} For three distinct involutions with commuting
endpoints and adjacent products of order at least three, IP is equivalent
to \(G_{01}\cap G_{12}=G_1\).

\emph{Proof.} Necessity is immediate. For sufficiency, the endpoint subgroup
is \(G_{02}=\{1,r_0,r_2,r_0r_2\}\). If \(G_{02}\cap G_{01}\) contains
\(r_2\) or \(r_0r_2\), then \(r_2\in G_{01}\cap G_{12}=G_1\), contradicting
distinctness. Thus this intersection is \(G_0\); dually
\(G_{02}\cap G_{12}=G_2\). These equalities exclude every generator from
the subgroup on the other two. Intersections involving singletons now
follow, and intersections involving the empty or full index set are
immediate. This covers all subset pairs. \(\square\)

\section{Deletion maps and supported extensions}\label{deletion-maps-and-supported-extensions}

A deletion for \(S\subseteq I\) is an endomorphism \(e_S:G\to G\)
satisfying \(e_S(r_i)=r_i\) for \(i\in S\) and \(e_S(r_i)=1\) otherwise.
Requiring these values matters: a general retraction onto \(G_S\) need
not be a deletion.

\textbf{Proposition UA.D (Boolean deletions).} If a marked group has a
deletion \(e_S\) for every \(S\subseteq I\), then \[
e_Se_T=e_{S\cap T},\qquad \operatorname{im}e_S=\operatorname{Fix}(e_S)=G_S.
\tag{D1}
\] The marking has IP, and \[
\ker e_S=\langle\!\langle r_i:i\notin S\rangle\!\rangle_G,
\qquad G=(\ker e_S)\rtimes G_S.
\tag{D2}
\]

\emph{Proof.} The two maps in the first equality agree on every
generator. The image is generated by precisely the retained markers, and
each of those markers is fixed; this proves the image and fixed-point
statements. An element of \(G_S\cap G_T\) is fixed by
\(e_Se_T=e_{S\cap T}\), so belongs to \(G_{S\cap T}\). The reverse
inclusion is immediate.

Let \(M\) be the normal closure in (D2). Clearly \(M\leq\ker e_S\). In
\(G/M\), every omitted marker is trivial, so the image of any word
equals the image of its deletion. A word in \(\ker e_S\) is therefore
trivial in \(G/M\), proving equality. Finally, for \(g\in G\), the
factorization \(g=(g e_S(g)^{-1})e_S(g)\) has factors in \(\ker e_S\)
and \(G_S\). Their intersection is trivial because \(e_S\) fixes
\(G_S\). This is the asserted semidirect product. \(\square\)

We now assemble extensions whose extra kernel is invisible on each
proper part of its support. The groups and kernels may be infinite and
nonabelian.

\textbf{Theorem UA.SK (supported-kernel realization).} Let \(I\) be
finite and let \(\mathcal F\) be a finite family of distinct subsets of
\(I\). Let \(B=\langle b_i:i\in I\rangle\) have IP. For every
\(S\in\mathcal F\), suppose the following data are given:

\begin{enumerate}
\def\labelenumi{\arabic{enumi}.}
\tightlist
\item
  A deletion homomorphism \(p_S:B\to B_S\), fixing \(b_i\) for
  \(i\in S\) and killing \(b_i\) otherwise.
\item
  A group \(E_S=\langle e_{S,i}:i\in S\rangle\) and a marked epimorphism
  \(q_S:E_S\to B_S\).
\item
  Injectivity of \(q_S\) on \((E_S)_T=\langle e_{S,i}:i\in T\rangle\)
  for every \(T\subsetneq S\).
\end{enumerate}

Put \(N_S=\ker q_S\), and form the group \[
H=\left\{(b,(x_S)_{S\in\mathcal F})\in
B\times\prod_{S\in\mathcal F}E_S:
q_S(x_S)=p_S(b)\text{ for every }S\right\}.
\tag{K1}
\] For \(i\in I\), define \(h_i\) to have base coordinate \(b_i\) and
\(S\)-coordinate \(e_{S,i}\) when \(i\in S\), and \(1\) otherwise. These
elements generate all of \(H\). If \(\pi:H\to B\) is the base projection
and \(H_L=\langle h_i:i\in L\rangle\), then, for every \(L\subseteq I\),
\[
\pi(H_L)=B_L,\qquad
\ker\pi=\prod_{S\in\mathcal F}N_S,\qquad
\ker\pi\cap H_L=\prod_{\substack{S\in\mathcal F\\S\subseteq L}}N_S.
\tag{K2}
\] The factors on the right are embedded as pure coordinate subgroups.
The marking of \(H\) has IP. If \(B\) and all \(E_S\) are finite, then
\[
|H_L|=|B_L|\prod_{\substack{S\in\mathcal F\\S\subseteq L}}|N_S|.
\tag{K3}
\]

\emph{Proof.} The compatibility equations in (K1) are preserved by
products and inverses, so \(H\) is a subgroup of the displayed direct
product. The prescribed tuple \(h_i\) satisfies those equations by the
literal deletion rule. Let \(H'\) be the subgroup generated by these
tuples. Its base image is \(B\), and the projection of
\(H'_S=\langle h_i:i\in S\rangle\) onto \(E_S\) is onto.

We first prove, by induction on inclusion among the finitely many
supports, that every pure coordinate element of \(N_S\) belongs to
\(H'_S\). Fix \(n\in N_S\). Choose a word on the markers indexed by
\(S\) which evaluates to \(n\) in \(E_S\), and evaluate the same word in
\(H'\), obtaining \(w\in H'_S\). Its base coordinate lies in \(B_S\) and
is \(q_S(n)=1\).

Consider its coordinate at another support \(T\). Because the word uses
only markers from \(S\), that coordinate lies in \((E_T)_{S\cap T}\).
Its image under \(q_T\) is \(p_T(1)=1\). If \(T\not\subseteq S\), then
\(S\cap T\subsetneq T\), and the proper-parabolic injectivity hypothesis
forces this coordinate to be \(1\). The only possible nonidentity
coordinates of \(w\) are consequently at supports \(T\subseteq S\). Its
\(S\)-coordinate is the prescribed \(n\), while every coordinate at a
proper support \(T\subsetneq S\) belongs to \(N_T\).

For each such proper support, the induction hypothesis supplies the pure
coordinate element having exactly that value in \(N_T\), as an element
of \(H'_T\leq H'_S\). Multiply \(w\) by the inverses of those elements.
Distinct pure coordinate subgroups commute in the ambient direct
product, so these operations remove the unwanted coordinates without
changing the \(S\)-coordinate. The result is the pure element \(n\) in
coordinate \(S\), and remains in \(H'_S\). A minimal support has no
smaller coordinates to remove, so this argument includes the initial
step. No commutativity within an individual kernel has been used.

The kernel of \(\pi\) in the full group \(H\) is exactly the direct
product of the \(N_S\), by (K1). The induction places that entire kernel
in \(H'\). Since \(H'\) already maps onto \(B\), every \(h\in H\) has
the same base coordinate as some \(h'\in H'\); then
\(h(h')^{-1}\in\ker\pi\leq H'\). Thus \(H'=H\).

The image equality \(\pi(H_L)=B_L\) follows on generators. If
\(x\in H_L\cap\ker\pi\), its \(S\)-coordinate lies in
\((E_S)_{L\cap S}\) and in \(N_S\). This coordinate is trivial whenever
\(S\not\subseteq L\), again by proper-parabolic injectivity. Conversely,
when \(S\subseteq L\), the isolation argument placed the entire pure
\(N_S\) in \(H_S\leq H_L\). This proves all three assertions in (K2).

The coordinate products in (K2) satisfy \[
\left(\prod_{S\subseteq L}N_S\right)\cap
\left(\prod_{S\subseteq M}N_S\right)
=\prod_{S\subseteq L\cap M}N_S.
\] Theorem UA.K now proves IP in \(H\). More explicitly, an element of
\(H_L\cap H_M\) has base in \(B_{L\cap M}\); subtracting a lift from
\(H_{L\cap M}\) leaves an element in both displayed kernel products,
hence in the kernel product for \(L\cap M\). This also shows exactly
where IP of the base is used. Finally (K3) is the order formula for the
surjection \(H_L\to B_L\) with the kernel just computed. \(\square\)

\textbf{Corollary UA.SK.Local (recognition on an antichain).} Under
UA.SK, if \(\mathcal F\) is an antichain under inclusion, then the
coordinate projection \(H_S\to E_S\) is a marked isomorphism for every
\(S\in\mathcal F\).

\emph{Proof.} It is onto on generators. If an element maps to \(1\), its
base coordinate also vanishes, since the base value of a word on \(S\)
is the \(q_S\)-image of its \(S\)-coordinate. Formula (K2) says that the
remaining kernel in \(H_S\) has only the factor \(N_S\): the antichain
excludes every other support contained in \(S\). The \(S\)-coordinate
detects that factor, so the element is \(1\). \(\square\)

\textbf{Corollary UA.SK.Delete (derived deletions).} If, in addition,
\(B\) has deletions for every subset of \(I\), then every \(E_S\) and
\(H\) has all its Boolean deletions.

\emph{Proof.} For \(T\subsetneq S\), the restriction
\(q_S:(E_S)_T\to B_T\) is a marked isomorphism: it is onto on generators
and injective by hypothesis. Compose \(q_S\), the deletion
\(B_S\to B_T\), and the inverse of this isomorphism. This gives a
homomorphism \(E_S\to(E_S)_T\) that fixes exactly the \(T\)-markers and
kills the others. For \(T=S\) use the identity. Proposition UA.D
supplies the Boolean identities.

For \(L\subseteq I\), apply the base deletion \(e_L\) and, in coordinate
\(S\), the local deletion onto \(S\cap L\). The two sides of each
compatibility equation agree on the local markers; the base deletions
commute by UA.D. Thus the coordinatewise operation preserves (K1). On
\(h_i\) it fixes \(h_i\) if \(i\in L\), and kills it otherwise. It is
the required deletion of \(H\). \(\square\)

\textbf{Corollary UA.SK.Central (lower-central kernel).} Under UA.SK,
suppose \(c\geq1\), \(\gamma_c(B)=1\), and \(N_S=\gamma_c(E_S)\) for
every support. Then \[
\gamma_c(H)=\ker\pi.
\tag{K4}
\]

\emph{Proof.} Applying \(\pi\) to a commutator of weight \(c\) gives
\(1\), so \(\gamma_c(H)\leq\ker\pi\). For the converse, repeat the
inclusion induction in the proof of UA.SK. The projection \(H_S\to E_S\)
is onto; therefore its image on \(\gamma_c(H_S)\) is
\(\gamma_c(E_S)=N_S\). Given \(n\in N_S\), choose a lift
\(w\in\gamma_c(H_S)\). Its base is \(1\), since \(\gamma_c(B)=1\). The
same proper-support argument makes all coordinates at
\(T\not\subseteq S\) trivial. The remaining proper-support coordinates
lie in \(N_T\); by induction their pure copies belong to
\(\gamma_c(H)\). Removing them leaves the pure \(n\), still in
\(\gamma_c(H)\). All pure kernels are obtained, proving (K4).
\(\square\)

\textbf{Corollary UA.SK.Separator (factorization across a separator).}
Under UA.SK, let \(I=L\cup M\). Suppose \(B=B_LB_M\) as a product of
sets, and every support in \(\mathcal F\) is contained in \(L\) or in
\(M\). Then \(H=H_LH_M\).

\emph{Proof.} For \(h\in H\), factor \(\pi(h)=b_Lb_M\) and choose lifts
\(u_L\in H_L\), \(u_M\in H_M\). Then \(h=n u_Lu_M\) for some
\(n\in\ker\pi\). Assign every pure coordinate factor of \(n\) to one of
the two sides containing its support. Formula (K2) gives \(n=n_Ln_M\),
with \(n_L\in H_L\), \(n_M\in H_M\). Each pure \(N_S\) is normal in
\(H\), because conjugation in the ambient product acts coordinatewise
and \(N_S\unlhd E_S\). Hence the product of the assigned \(M\)-factors
is normal in \(H\). Moving \(n_M\) past \(u_L\) gives \[
h=(n_Lu_L)(u_L^{-1}n_Mu_L\,u_M)\in H_LH_M.
\] The opposite inclusion is automatic. \(\square\)

The hypotheses of UA.SK do not identify \(H\) with a universal amalgam
of the \(E_S\). Additional cross-relations and a
presentation-recognition argument may be necessary. Nor can injectivity
on every proper parabolic be replaced by injectivity on individual
markers alone.

\textbf{Proposition UA.Fiber (two compatible split factors).} Let \(A\)
and \(B\) be marked groups on finite label sets \(L\) and \(M\), and put
\(J=L\cap M\). Suppose their \(J\)-parabolics are identified with the
same marked group \(C\), and that \(p_A:A\to C\) and \(p_B:B\to C\) are
deletions. Then \[
A\times_C B\cong(\ker p_A\times\ker p_B)\rtimes C.
\tag{F1}
\] This full fiber product is generated by the tuples \((a_i,1)\) for
\(i\in L\setminus J\), \((1,b_i)\) for \(i\in M\setminus J\), and
\((c_i,c_i)\) for \(i\in J\). The two original groups embed through
\(a\mapsto(a,p_A(a))\) and \(b\mapsto(p_B(b),b)\). If finite, its order
is \(|A||B|/|C|\).

\emph{Proof.} Every compatible pair has a unique form \((n_Ac,n_Bc)\),
giving the split decomposition and the order formula. The stated copies
of \(A\) and \(B\) are generated by subsets of the stated tuple marking.
They contain \((n_A,1)\) and \((1,n_B)\), respectively, while either
contains the diagonal copy of \(C\). These subgroups generate the split
product, proving the assertion. The embedding formulas are injective
because their first and second projections, respectively, are
identities. \(\square\)

\section{Recognition of presentations and compatible products}

All homomorphisms in this paper are group homomorphisms. A product of subgroups
means a product of sets until it has been shown to be a subgroup. No group is
assumed finite in this section.

\textbf{Theorem UA.Normal (normal kernel recognition).}
Let $\pi:P\to B$ and $\rho_i:P\to E_i$, for $i$ in a finite set $J$, be
homomorphisms. Suppose $K_i\unlhd P$ satisfy
\begin{align}
 \ker\pi&=\gen{K_i:i\in J},\label{eq:basekernel}\\
 \rho_i(K_j)&=1\quad(i\ne j),\label{eq:offdiag}\\
 \ker\rho_i\cap K_i&=1.\label{eq:faithful}
\end{align}
Then the homomorphism
\[
 P\longrightarrow B\times\prod_{i\in J}E_i,
 \qquad x\longmapsto (\pi(x),(\rho_i(x))_i)
\]
is injective. Moreover, the kernel in \eqref{eq:basekernel} is the internal
direct product of the $K_i$. For injectivity alone it suffices to replace
\eqref{eq:basekernel} by $\ker\pi\leq\gen{K_i:i\in J}$.

\begin{proof}
For $i\ne j$, an element of $K_i\cap K_j$ is killed by $\rho_i$ by
\eqref{eq:offdiag}, and is therefore trivial by \eqref{eq:faithful}.
Normality gives $[K_i,K_j]\leq K_i\cap K_j$, so distinct factors commute.
Their finite set product is consequently their generated subgroup. If
$\prod_i k_i=1$ with $k_i\in K_i$, applying $\rho_j$ gives
$\rho_j(k_j)=1$, whence $k_j=1$. Thus multiplication identifies the external
product with the generated subgroup. Finally, an element in the kernel of the
displayed homomorphism belongs to $\ker\pi$, hence has such a product
expression. Every coordinate vanishes, so every factor is trivial.
\end{proof}

\textbf{Corollary UA.RecognizeQuotient.}
Let $f:P\twoheadrightarrow H$ be a homomorphism, and let $H$ have a base
homomorphism $p:H\to B$ and coordinates $c_i:H\to E_i$. If
UA.Normal applies to $p f$, $c_i f$ and normal subgroups
$K_i\leq P$, then $f$ is an isomorphism.

\begin{proof}
The coordinate map factors through $f$, so injectivity of the coordinate map
implies injectivity of $f$. Surjectivity is the stated hypothesis.
\end{proof}

Normality is substantial. Let $P=D_8=\gen{u,v\mid u^2=v^2=(uv)^4=1}$,
let $B=1$, and take $K_1=\gen u$, $K_2=\gen v$. Its abelianization
$P\to C_2^2$ has coordinates faithfully detecting the two factors and killing
the other one. The factors generate $P$, but the map has the nontrivial kernel
$\gen{(uv)^2}$. Here neither factor is normal. Thus a correct local model and
coordinate detection do not themselves eliminate mixed relations.

The following specialization uses the full compatible product $H$ in
UA.SK and assumes that its support family is an antichain.

\textbf{Theorem UA.Recognition (complete presentation).}
In this setting, let $P=\gen{r_t:t\in I}$ be given by a
presentation. Suppose:
\begin{enumerate}
\item its relators hold on the marker tuples in the compatible product $H$;
\item the local relations of every $E_S$ hold on $(r_t)_{t\in S}$;
\item the image $K_S$ of $N_S$ under $E_S\to P$ is normal in $P$;
\item adjoining the relations $K_S=1$ for all $S$ gives exactly the marked
presentation of $B$.
\end{enumerate}
Then the marker map $P\to H$ is an isomorphism. In particular the $K_S$
form an internal direct product, all marked intersections are exact, and if
the input groups are finite then
\begin{equation}
 |P_L|=|B_L|\prod_{S\subseteq L}|N_S|.                         \label{eq:orders}
\end{equation}
The additional deletion conclusion of UA.SK.Delete transfers to $P$ under
its hypotheses. The same holds for the geometric conclusion of
UA.SK.Geometry under its hypotheses.

\begin{proof}
The first assumption and the generation assertion of UA.SK give a marked
surjection $f:P\twoheadrightarrow H$. The second gives $E_S\to P_S$;
its composite with $f$ and projection to $E_S$ is the identity on every
local marker, hence is the identity homomorphism. Thus each local map is
injective, without an order comparison. Under $f$, the image of $N_S$ is
its pure $S$-coordinate: the antichain makes every other intersection of
supports proper, where $q_T$ is injective. Consequently the coordinate maps
kill $K_S$ outside its own coordinate and are injective on $K_S$ in that
coordinate.

Normality makes $K=\gen{K_S:S\in\mathcal S}$ a normal subgroup. The last
assumption identifies the marked map $P/K\to B$ with an isomorphism, so
the base map $P\to B$ has kernel exactly $K$. UA.RecognizeQuotient
now proves that $f$ is an isomorphism. The remaining conclusions follow
from UA.SK and the explicitly indicated additional clauses.
\end{proof}

The last assumption is a presentation calculation. It does not assume the
isomorphism $P\cong H$ which is being proved. In practice it says that, after
killing the displayed local kernel words, the remaining relations are the
known base presentation. A pair of marked maps which are inverse on generators
also proves this condition. Normality likewise has a concrete local test: the
markers within $S$ already normalize $K_S$, and one must check the remaining
markers. It is enough that those remaining markers centralize a generating
set of $K_S$. Merely writing down the desired quotient does not supply this
check.


\textbf{Proposition UA.Mix (jointly faithful deletions).}
Let $G=\gen{r_i:i\in I}$ have every Boolean deletion $e_J$. Suppose
$I=L\cup M$, put $C=L\cap M$, and assume
$\ker e_L\cap\ker e_M=1$. Then
\[
 G\longrightarrow G_L\times_{G_C}G_M,\qquad
 g\longmapsto(e_L(g),e_M(g))
\]
is a marked isomorphism. The maps to $G_C$ on the right are the restrictions
of the corresponding deletions. No group is assumed finite.
\begin{proof}
UA.D gives $e_Le_M=e_C=e_Me_L$. Hence the displayed pair is compatible,
and the homomorphism's kernel is the stated trivial intersection.
For a compatible pair $(x,y)\in G_L\times_{G_C}G_M$, put
$c=e_C(x)=e_C(y)$. Since the two deletions fix their respective
parabolics, $e_M(x)=c=e_L(y)$ and both deletions fix $c$. Thus
\[
 e_L(xc^{-1}y)=xc^{-1}c=x,\qquad
 e_M(xc^{-1}y)=cc^{-1}y=y.
\]
Every compatible pair therefore has the explicit preimage $xc^{-1}y$.
The homomorphism is surjective as well as injective, with inverse given
by this formula. Its values on the markers give the stated marking.
\end{proof}

The kernel condition can be checked by coordinates. In a supported-kernel
realization with inherited deletions, suppose every base coordinate is
detected on at least one side, and every support is contained in $L$ or
$M$. An element killed by both deletions has trivial base coordinate. Each
pure supported-kernel coordinate is then detected on a side containing its
support, so all these coordinates are trivial too. This proves the kernel
condition in the path and two-branch applications below; it is not assumed
from a coincidence of finite orders.

\section{Coordinate kernels and incidence geometries}\label{coordinate-kernels-and-incidence-geometries}

An incidence system has a finite type set, objects carrying types, and a
symmetric incidence relation; distinct objects of one type are not
incident. A flag is a set of pairwise incident objects. A chamber has
one object of each type. A geometry is an incidence system in which
every flag extends to a chamber. The residue of a flag consists of the
objects of the remaining types incident with that flag, with induced
incidence. The geometry is residually connected when the incidence graph
of every residue of rank at least two, including the empty-flag residue,
is connected. It is thin when every panel, a flag missing just one type,
lies in exactly two chambers. A regular hypertope is a thin, residually
connected geometry admitting a chamber-transitive type-preserving
automorphism group.

For a marked group indexed by \(I\), its maximal parabolics are
\(P_i=G_{I\setminus\{i\}}\). The associated typed coset incidence system
has objects \(P_i g\), and incidence means nonempty intersection of
cosets. All actions and cosets below use this right-coset convention.

For subgroups \(P_i\leq G\), typed coset Helly means that every pairwise
intersecting family of right cosets \(P_i g_i\), with distinct indices,
has a common element. This property is separate from IP. It supplies a
common representative for an incident flag, whereas IP determines the
stabilizer of a flag already represented in this way.

\textbf{Proposition UA.Geometry (coset recognition and faithful chamber
action).} Suppose the marking of \(G\) consists of nonidentity
involutions, has IP, and has typed coset Helly for its maximal
parabolics. The coset incidence system is then a thin, residually
connected geometry, and \(G\) is its full type-preserving automorphism
group, acting regularly on chambers. More generally, the full
type-preserving automorphism group of any thin, residually connected
geometry of finite rank acts freely on chambers. A chamber-transitive
subgroup is consequently the full group.

\emph{Proof.} In the coset system, a common coset point extends every
flag to a chamber by choosing the coset of each missing type through
that point. Right translation sends every such flag to its base flag.
For a base flag of type set \(S\), IP gives stabilizer
\(G_{I\setminus S}\). In particular the full chamber stabilizer is
\(1\), while a panel stabilizer is a group generated by one nonidentity
involution and has order two. There are therefore exactly two chambers
through each panel.

Consider a residue with remaining type set \(T\), where \(|T|\geq2\).
Its stabilizer is \(G_T\), and its maximal object stabilizers are
\(G_{T\setminus\{i\}}\), \(i\in T\). Every marker indexed by \(T\) lies
in at least one of these proper parabolics. A word in those subgroups
supplies a walk between their translated base objects in the residue
incidence graph; all objects are translates of base objects and all base
objects lie in a common flag. The residue is therefore connected.

For the general assertion, first observe that every residually connected
finite-rank geometry has a connected chamber graph. This follows by
induction on rank: chambers through a fixed object are connected through
galleries in its residue, and a path in the incidence graph, extending
each consecutive incident pair to a chamber, connects arbitrary chambers
through these subgraphs. The rank-one case holds because all chambers
share the empty panel. In a thin geometry an automorphism fixing a
chamber fixes its unique neighbor across each panel. It therefore fixes
every chamber along every gallery, and hence every object. The chamber
action is free. If a subgroup is chamber-transitive, any full
automorphism agrees on a base chamber with one subgroup element; their
quotient fixes that chamber and is the identity. Apply this to the coset
geometry to obtain the stated full automorphism group. \(\square\)

\subsection{Helly transfer through coordinate
kernels}\label{helly-transfer-through-coordinate-kernels}

Write right cosets as \(P g\), with the subgroup on the left. Let
\(q:H\to B\) be an epimorphism of groups generated by corresponding
marked families \((h_i)_{i\in I}\) and \((b_i)_{i\in I}\), where \(I\)
is finite. For \(J\subseteq I\), write \(H_J=\langle h_j:j\in J\rangle\)
and \(B_J=\langle b_j:j\in J\rangle\). In particular, \(q(H_J)=B_J\) for
every J.

Fix any set \(\mathcal J\) of subsets of I as the permitted types. The
maximal-parabolic case is \(\mathcal J=\{I\setminus\{i\}:i\in I\}\).

\textbf{Theorem UA.SK.Helly (coordinate-kernel transfer).} Suppose B has
IP. Suppose \(N=\ker q\) is an internal direct product
\(\prod_{\alpha\in A}N_\alpha\), with a finite factor set A, and for
every \(J\in\mathcal J\), \[
 N\cap H_J=\prod_{\alpha\in A_J}N_\alpha
 \qquad\text{for some }A_J\subseteq A.
 \tag{H1}
\] If the family \((B_J)_{J\in\mathcal J}\) has typed coset Helly, then
so does \((H_J)_{J\in\mathcal J}\).

No commutativity within the factors, finiteness of H or B, or
faithfulness of an action of B on N is assumed. The required
intersection-image fact is \(q(H_{J\cap K})=B_J\cap B_K\); it follows
here from the marked map and IP of B. IP of H is not an additional
assumption for this transfer theorem.

\emph{Proof.} Consider pairwise intersecting cosets \(H_{J_i}g_i\).
Their images \(B_{J_i}q(g_i)\) intersect pairwise. Choose a common base
element b by the assumed Helly property, and choose u in H with
\(q(u)=b\). Each set \[
 F_i=H_{J_i}g_i\cap q^{-1}(b)
\] is nonempty: because \(bq(g_i)^{-1}\in B_{J_i}\), a preimage of that
element in \(H_{J_i}\), multiplied on the right by \(g_i\), belongs to
\(F_i\).

We first show that \(F_i\cap F_j\ne\varnothing\). Choose
\(x\in H_{J_i}g_i\cap H_{J_j}g_j\) and put z=q(x). Both b and z belong
to each of the two base cosets, so \[
 bz^{-1}\in B_{J_i}\cap B_{J_j}=B_{J_i\cap J_j}.
\] The marked map is surjective on this intersection parabolic. Choose
\(h\in H_{J_i\cap J_j}\) with \(q(h)=bz^{-1}\). Since h belongs to both
\(H_{J_i}\) and \(H_{J_j}\), left multiplication by h preserves both
right cosets. Thus hx belongs to their intersection and has image b.
This proves the assertion without assuming that q carries an arbitrary
intersection onto the intersection of its images.

Choose \(x_i\in F_i\) and put \(k_i=x_i u^{-1}\in N\). If \(y\in F_i\),
then \(yx_i^{-1}\in N\cap H_{J_i}\); the converse also holds. Hence \[
 F_i u^{-1}=(N\cap H_{J_i})k_i.
 \tag{H2}
\] Under the direct product coordinates of N, (H1) makes (H2)
unrestricted in coordinates \(\alpha\in A_{J_i}\), and fixes coordinate
\(\alpha\notin A_{J_i}\) to \((k_i)_\alpha\). Pairwise intersections
imply that any two constraints fixing the same coordinate fix it to the
same value. Choose that value in every constrained coordinate and the
identity in every unconstrained coordinate. The resulting element
belongs to all sets in (H2). Multiplication on the right by u supplies a
common element of the original cosets. This proves the theorem.
\(\square\)

\textbf{Corollary UA.SK.Geometry.} In the supported-kernel realization
UA.SK, assume the base marking gives a regular hypertope and every
marked generator in every attached group is an involution. Then the
resulting marking also gives a regular hypertope, and its full
type-preserving automorphism group is the constructed group.

\emph{Proof.} UA.SK gives IP and the exact profile \[
 N\cap H_J=\prod_{S\subseteq J}N_S.
\] The base maximal-parabolic cosets have typed Helly. Theorem
UA.SK.Helly applies with factors indexed by the supports S. Each
constructed generator is a tuple whose base coordinate and every
nontrivial attached coordinate are involutions; hence its square is the
identity. Its base image is nonidentity, so its order is exactly two.
UA.Geometry now gives thinness, residual connectedness, chamber
regularity and equality with the full type-preserving automorphism
group. In particular, this conclusion uses typed Helly in addition to
IP. \(\square\)

A sufficient alternative to local involutivity is that every support has
size at least two: each singleton local parabolic is then proper, and
injectivity into its nonidentity involutory base image forces the local
marker to have order two. This is a sufficient condition, not an
equivalent reformulation. Without either condition the geometric
corollary fails. With a singleton type set and support, the quotient
\(C_4\to C_2\) is injective on its only proper parabolic, the trivial
subgroup. UA.SK then constructs \(C_4\), whose distinguished generator
has order four. In rank three, adjoining two independent \(C_2\) markers
gives \(C_4\times C_2^2\) with the same defective singleton attachment
and a four-element panel.

\section{Frattini coordinates}\label{frattini-coordinates}

\textbf{Theorem UA.Phi (primitive saturation and coordinates).} Let
\(P\) be a finite \(p\)-group with an irredundant generating marking
\((r_i)_{i\in I}\), \(|I|=d\), and let \(F=\Phi(P)\). Then, for every
\(S\subseteq I\), \[
\Phi(P_S)=P_S\cap F,\qquad
P_SF=\pi^{-1}\operatorname{span}(\bar r_i:i\in S),\qquad
[P:P_SF]=p^{d-|S|}.
\tag{P1}
\] The saturated parabolics \(P_SF\) intersect according to intersection
of index sets. Each omission subgroup \(P_{I\setminus\{i\}}\) has
exactly one maximal overgroup, namely \(P_{I\setminus\{i\}}F\).

If \(B\geq F\) and \((s_i)_{i\in J}\) is an irredundant relative
generating family over \(B\), then \(B\unlhd P\), \(P/B\) has basis
\((s_iB)_{i\in J}\) over \(\mathbb F_p\), and \[
\langle B,s_i:i\in T\rangle/B=\operatorname{span}(s_iB:i\in T).
\tag{P2}
\] If \(p=2\) and all original \(r_i\) are involutions, then \(F=P'\)
and \(\Phi(P_S)=P_S'\) for every \(S\).

\emph{Proof.} Burnside's basis theorem makes the \(\bar r_i\) a basis of
\(P/F\). The image of \(P_S\) therefore has order \(p^{|S|}\), with
kernel \(P_S\cap F\). Since this image is elementary abelian,
\(\Phi(P_S)\leq P_S\cap F\). On the other hand \(P_S\) is generated by
\(|S|\) elements, so Burnside gives \([P_S:\Phi(P_S)]\leq p^{|S|}\). The
quotient by \(P_S\cap F\) already has that order. Equality follows,
proving the first formula. The remaining formulas and the intersection
assertion are inverse-image statements for coordinate subspaces.

Every maximal subgroup of a finite \(p\)-group contains \(F\) and
corresponds to a hyperplane of \(P/F\). A hyperplane containing a
specified coordinate hyperplane equals it, proving uniqueness of the
maximal overgroup of an omission subgroup.

For the relative assertion, every subgroup containing \(F\) is normal,
because \(P/F\) is abelian. The quotient \(P/B\) is elementary abelian.
The indicated images span it, and any nontrivial linear relation would
express one \(s_i\) using \(B\) and the other markers, contradicting
irredundancy. They form a basis, proving (P2). Finally, an abelian group
generated by involutions is elementary abelian of exponent two. Thus
\(P'\leq\Phi(P)\), while the elementary abelian quotient \(P/P'\) gives
the reverse inclusion. The same argument applies to each \(P_S\).
\(\square\)

These coordinates also describe the corresponding omission cosets: a
coset prescribes one coordinate value, and every simultaneous assignment
gives a chamber. The geometry is complete multipartite, with \(p\)
objects in each type, and the quotient translation group acts regularly
on chambers. For \(p>2\) it need not be the full type-preserving
automorphism group; the thin case \(p=2\) is the one covered by
UA.Geometry.

\section{The canonical quotient of a forest marking}\label{the-canonical-quotient-of-a-forest-marking}

Suppose a group is generated by nonidentity involutions \(r_i\),
\(i\in I\). Its noncommutation graph has vertex set \(I\), and an edge
\(\{i,j\}\) precisely when \(r_i\) and \(r_j\) do not commute. We write
\(G_J=\langle r_j:j\in J\rangle\), as in the marking definitions. The intersection
property in this section always refers to the original group. It is not
assumed to survive an intermediate quotient.

\textbf{Lemma UA.Lift (descending lift).} Let \(G\) be a finite
nilpotent group, let \(A,B\le G\) satisfy \(A\cap B=1\), and fix an
integer \(m\ge1\). Put \(D_j=\gamma_j(G)\), \(A_j=A\cap D_j\),
\(B_j=B\cap D_j\). If the images of \(A_j,B_j\) generate \(D_j/D_{j+1}\)
for every \(j\ge m\), then \(D_j=A_jB_j\) for every \(j\ge m\).

\textbf{Proof.} Descend from a trivial term. Assuming
\(D_{j+1}=A_{j+1}B_{j+1}\), if \(a\in A_j,b\in B_j\) have equal images
in the next quotient, write \(ab^{-1}=a'b'\) with
\(a'\in A_{j+1},b'\in B_{j+1}\). Then \(a'^{-1}a=b'b\in A\cap B\), so
both images are trivial. Their images therefore form an internal direct
sum in the abelian quotient. Counting gives
\(|D_j|=|A_j||B_j|=|A_jB_j|\), proving equality as sets. No prior
normality or permutability of \(A,B\) is assumed. ∎

\textbf{Lemma UA.Support (graded support).} Let a positively graded Lie
algebra be generated in degree one by disjoint sets \(X,Y\), with
subalgebras \(L_X,L_Y\). If \([x,y]\in L_Y\) for \(x\in X,y\in Y\), then
\(L=L_X+L_Y\), degree by degree.

\textbf{Proof.} The derivation \(\operatorname{ad}x\) preserves \(L_Y\)
because it does so on its generators. Iterated brackets of those
derivations also preserve \(L_Y\), by Jacobi. Hence
\([L_X,L_Y]\subseteq L_Y\), and \(L_X+L_Y\) is a Lie subalgebra
containing all generators. The argument holds in characteristic two. ∎

\textbf{Theorem UA.Forest (forest quotient).} Suppose \(G\) is a finite
two-group generated by nonidentity involutions \(r_i\), with marked IP,
and its noncommuting pairs form a forest \(T\) on \(d\) vertices and
\(m\) edges. Then edge commutators are independent in
\(G'/\gamma_3(G)\), and \(G/\gamma_3(G)\) is the marked group \[
 (v,z)(w,t)=(v+w,z+t+\beta(v,w)),\quad
 v,w\in\mathbb F_2^d,\ z,t\in\mathbb F_2^m,               \tag{F0}
\] where, after orienting each edge, its component of \(\beta\) is
\(v_{\rm head}w_{\rm tail}\). Its order is \(2^{d+m}\).

\textbf{Proof.} IP gives irredundancy. UA.Phi gives \(G/G'=C_2^d\) and
\(G'\cap\langle r_i\rangle=1\). The lower-central graded Lie ring is
generated by these degree-one elements; all its graded components are
\(\mathbb F_2\)-spaces by bilinearity. Edge commutators span degree two.
If they satisfy a nonzero linear relation, root each forest component
and choose a participating edge whose child endpoint has maximal depth
in its component. Removing it separates the descendant set \(X\) from
all other vertices \(Y\), and no other participating edge is internal to
\(X\). The relation therefore writes the only nonzero cross bracket as a
sum of brackets internal to \(Y\). Lemma UA.Support shows every graded
component is generated by pure \(X\)- and pure \(Y\)-brackets. Their
representatives lie in \(G_X\cap\gamma_j(G)\) and
\(G_Y\cap\gamma_j(G)\).

Actual IP gives \(G_X\cap G_Y=1\), so Lemma UA.Lift gives
\(G'=(G'\cap G_X)(G'\cap G_Y)\). For the crossing endpoints
\(x\in X,y\in Y\), write \([r_x,r_y]=uv\) accordingly. Then
\(u\in G_X\cap G_{\{x\}\cup Y}=\langle r_x\rangle\), whence \(u=1\);
next \([r_x,r_y]\in G_Y\cap G_{\{x,y\}}=\langle r_y\rangle\), so it too
is one. This contradicts that the edge is noncommuting. Independence
follows.

Bilinearity makes (F0) associative; its coordinate markers are
involutions with independent central edge commutators. It is generated
by those markers, and its parabolic on any subset \(J\) consists exactly
of vertex coordinates supported on \(J\) and edge coordinates supported
on edges internal to \(J\). Thus it has IP. Restricting the vertex
coordinates to a subset and retaining exactly its internal edge
coordinates preserves the multiplication formula, so it also gives every
Boolean deletion. Collection in the corresponding class-two presentation
gives at most \(2^{d+m}\) elements, and this coordinate model attains
the bound. Independence in \(G'/\gamma_3(G)\) gives the same order for
\(G/\gamma_3(G)\); hence the marked map is an isomorphism. ∎

\section{Finite quotients of prescribed excess order}\label{finite-quotients-of-prescribed-excess-order}

The next theorem separates a universal finite search space from the task
of proving that a particular family has the required class-two quotient.
That identification is an application hypothesis, not a consequence of
the search.

\textbf{Theorem UA.Nilpotent (complete marked kernel parameter space).}
Let \(F=\langle x_i:i\in I\rangle\), where \(|I|=d\) and \(x_i^2=1\).
Suppose \(F\) has all Boolean deletions. Put \(Q=F/\gamma_3(F)\), with
its image marking. For \(k\geq0\), set \[
U=F/\gamma_{k+3}(F),\qquad K=\gamma_3(U),\qquad
K_S=K\cap U_S.
\tag{N1}
\] Then \(U\) is a finite two-group, \[
|U|\leq 2^{\,\sum_{j=1}^{k+2}d^j},
\tag{N2}
\] the markings of \(U\) and \(Q\) have Boolean deletions and IP, and
\(U/K\) is canonically the marked group \(Q\). For every
\(S\subseteq I\), \[
K_S=\gamma_3(U_S).
\tag{N3}
\]

There is a bijection between the following two sets:

\begin{enumerate}
\def\labelenumi{\arabic{enumi}.}
\tightlist
\item
  Ordered marked isomorphism classes of finite two-group quotients
  \(f:F\twoheadrightarrow G\) for which the induced marked map
  \(Q\to G/\gamma_3(G)\) is an isomorphism, \(|G|=|Q|2^k\), and \(G\)
  has IP.
\item
  Normal subgroups \(N\unlhd U\) satisfying \(N\leq K\), \([K:N]=2^k\),
  and \[
  \frac{NK_S}{N}\cap\frac{NK_T}{N}
  =\frac{NK_{S\cap T}}{N}
  \quad(S,T\subseteq I).
  \tag{N4}
  \]
\end{enumerate}

The correspondence sends \(N\) to \(U/N\) with the original image
markers. If \(D\leq\operatorname{Sym}(I)\) is specified and each
\(\sigma\in D\) induces an automorphism \(x_i\mapsto x_{\sigma(i)}\) of
\(F\), then equivalence allowing precisely these label permutations is
the \(D\)-orbit relation on the kernels in item 2.

\emph{Proof.} The degree-one part of the lower-central graded Lie ring
of \(F\) is generated by elements of exponent dividing two. By
multilinearity, every iterated commutator in those generators also has
exponent dividing two in its homogeneous component. In degree \(j\), at
most \(d^j\) left-normed commutators suffice to generate that component.
Consequently \[
|\gamma_j(F)/\gamma_{j+1}(F)|\leq2^{d^j}.
\] Multiplying these bounds for \(1\leq j\leq k+2\) proves finiteness,
the two-power order and (N2). Each deletion preserves every
lower-central term, and therefore descends both to \(U\) and to \(Q\).
Proposition UA.D gives IP for both image markings; no separate IP
hypothesis on \(Q\) is needed. The epimorphic-image property of the
lower central series gives \(U/K\cong F/\gamma_3(F)=Q\).

For (N3), one inclusion follows from \(U_S\leq U\). If
\(z\in\gamma_3(U)\cap U_S\), the deletion \(e_S\) fixes \(z\) and sends
\(\gamma_3(U)\) into \(\gamma_3(U_S)\). This proves the reverse
inclusion.

Let \(G\) satisfy item 1. Its subgroup \(\gamma_3(G)\) has order
\(2^k\). In a nilpotent group, equality between two successive
nontrivial lower-central terms would persist at every later term,
contradicting termination of that series. Thus each nontrivial step
beginning at \(\gamma_3(G)\) strictly decreases its order. There are at
most \(k\) such steps, and \(\gamma_{k+3}(G)=1\). The map \(F\to G\)
therefore factors through \(U\). Its kernel \(N\) lies in \(K\), since
the induced map \(U/K\to G/\gamma_3(G)\) is an isomorphism, and the
orders give \([K:N]=2^k\).

Conversely, let \(N\unlhd U\) with \(N\leq K\) and \([K:N]=2^k\), and
write \(G=U/N\). Its third lower-central subgroup is \(K/N\), its
canonical class-two quotient is \(Q\), and its order is \(|Q|2^k\). If
\(q:U\to G\), then \[
(K/N)\cap G_S=q(K\cap U_S)=NK_S/N.
\tag{N5}
\] Indeed, \(q(u)\in K/N\) with \(u\in U_S\) implies \(u\in K\), because
\(q^{-1}(K/N)=K\). The reverse inclusion is immediate. Applying UA.K to
\(G\to Q\) proves that IP in \(G\) is equivalent to precisely (N4).

An isomorphism preserving each ordered marker commutes with the two maps
from \(U\), because it agrees with them on a generating family. Their
kernels are consequently equal. Equal kernels plainly give the same
marked quotient. Finally, every specified label permutation acts on
\(F\), preserves its characteristic lower-central terms, and hence acts
on \(U\) and on the displayed kernels. A marker-permuting isomorphism
induces the corresponding kernel equality after this automorphism, and
the converse descends that automorphism to the quotients. This proves
the orbit assertion at exactly the specified equivalence scope.
\(\square\)

\textbf{Theorem UA.Ladder (invariant index-two descent).} Let \(U\) be a
finite two-group and \(M\unlhd U\). The \(U\)-normal subgroups
\(L\leq M\) of index two are exactly the inverse images of the
hyperplanes of \[
M/\bigl(\Phi(M)[M,U]\bigr).
\tag{N6}
\] For \(K\unlhd U\), every \(U\)-normal subgroup \(N\leq K\) of index
\(2^k\) is obtained by \(k\) successive such steps.

\emph{Proof.} An index-two subgroup \(L\) of \(M\) contains \(\Phi(M)\)
and has quotient \(C_2\). If \(L\) is \(U\)-normal, conjugation by \(U\)
on \(M/L\) is trivial, since \(\operatorname{Aut}(C_2)=1\). Thus
\([M,U]\leq L\). Conversely, a hyperplane inverse image has index two
and contains \([M,U]\); for \(l\in L\) and \(u\in U\), the conjugate of
\(l\) differs from \(l\) by an element of \([M,U]\), so lies in \(L\).
The quotient in (N6) is elementary abelian by Burnside's basis theorem,
proving the first assertion.

Suppose \(N\unlhd U\) and \(N<M\). The nontrivial finite two-group
\(M/N\) has a nonzero Frattini quotient \(V\). Conjugation by \(U/N\)
acts on the nonzero linear maps \(V\to\mathbb F_2\). If \(\dim V=a\),
there are \(2^a-1\) such maps, an odd number. Every orbit has two-power
size, so at least one orbit has size one. The kernel of a fixed nonzero
map lifts to a \(U\)-normal subgroup between \(N\) and \(M\), of index
two in \(M\). Repeating with the new subgroup gives a chain terminating
at \(N\); its length is \(k\), since the total index is \(2^k\).
\(\square\)

\textbf{Lemma UA.Order (presentation recognition by a sharp bound).} Let
\(A_1,\ldots,A_t\) be finite normal subgroups of \(G\), with
\(|A_j|\leq a_j\). Suppose \([G:A_1\cdots A_t]\leq b\). Then
\(|G|\leq b\prod_j a_j\). If a finite quotient of \(G\) has this latter
order, its quotient map is an isomorphism.

\emph{Proof.} Normality makes each partial product a subgroup. The
multiplication map from \(A_1\times\cdots\times A_t\) onto the product
is surjective as a map of sets, so that product has order at most
\(\prod_j a_j\). Multiplication by the quotient index proves the bound
and finiteness of \(G\). A quotient attaining the bound forces equality
of the finite orders and hence a trivial kernel. \(\square\)

\textbf{Corollary UA.Search (rank and excess classify the finite
search).} Let T be a specified forest on d vertices and m edges, and let
k≥0. Set

\[
 F_T=\langle x_i\mid x_i^2=1,\ [x_i,x_j]=1\text{ for nonedges}\rangle,
 \quad U_{T,k}=F_T/\gamma_{k+3}(F_T),\quad K=\gamma_3(U_{T,k}).
\]

The ordered IP markings by nonidentity involutions with graph T and
order \(2^{d+m+k}\) are exactly the quotients \(U_{T,k}/N\), where N is
normal, contained in K, \([K:N]=2^k\), and the images
\((N(K\cap(U_{T,k})_J))/N\) have IP for every pair J,L. Generator-set
equivalence is exactly the action of Aut(T) on these kernels.

\emph{Proof.} Killing any set of generators preserves every defining
relator of \(F_T\). Collection in its class-two quotient gives at most
d+m binary coordinates, and (F0) realizes all of them, so this quotient
is precisely the marked model \(Q_T\). Every target has that same
quotient by UA.Forest. These are exactly UA.Nilpotent's hypotheses. Its
all-parabolic condition and UA.Ladder give the stated bijection and
exhaustive search. An ordered isomorphism must commute with the
generating map, hence has the same kernel. A permutation of the
generator set preserves the noncommutation graph and extends to an
automorphism of \(F_T\) and \(U_{T,k}\); this proves the last claim. ∎

The last corollary is a finite classification procedure, with the
marking retained. First form the finite group \(U_{T,k}\); then iterate
UA.Ladder inside \(\gamma_3(U_{T,k})\), test the explicit intersection
condition of UA.Nilpotent, and take the declared graph-automorphism
orbits. The theorem proves that this procedure terminates and misses no
target marking. Any reported numbers of groups, named presentations or
universal orders still require execution and an application-specific
identification. No numerical census is assumed or reported in this
paper.

\section{Affine square groups}\label{affine-square-groups}

Let \[
 a(x,y)=(-x,y),\quad b(x,y)=(y,x),\quad c(x,y)=(x,1-y).
\] \textbf{Proposition UA.Affine (the complete square presentation).}
These generate \(W=\mathbb Z^2\rtimes D_8\), of type \(\{4,4\}\), with
\(ac=ca\). Its presentation is \[
 W=\langle a,b,c\mid a^2=b^2=c^2=(ab)^4=(bc)^4=(ac)^2=1\rangle. \tag{S1}
\] \emph{Proof.} Put \(d=bab\), \(y_0=cd\), and \(x_0=by_0b\). The
relations give \([a,d]=1\), \([a,y_0]=1\), \(dy_0d=y_0^{-1}\),
\(ax_0a=x_0^{-1}\), and conjugation by \(b\) swaps \(x_0,y_0\).
Substitution of \(bc=x_0ab\) makes \((bc)^4=1\) equivalent to
\([x_0,y_0]=1\). Thus every word collects to \(x_0^u y_0^v M\),
\(M\in\langle a,b\rangle=D_8\). In the affine model \(x_0,y_0\) are the
unit translations, so these forms are distinct. This proves (S1).
\(\square\)

The normal form also fixes our affine multiplication convention: the
rightmost transformation acts first. The symbol \(D_8\) denotes the
group of signed permutation matrices in two coordinates, of order eight.
In all dihedral notation below the subscript denotes the order of the
group.

\textbf{Proposition UA.Square (finite square seeds).} For \(n\ge5\),
define the invariant lattice \[
 \Lambda_n=\begin{cases}
  2^e\mathbb Z^2,&n=2e+3,\\
  2^e\langle(1,1),(1,-1)\rangle,&n=2e+4,
 \end{cases}\qquad e\ge1,
\] and put \(H_n=W/\Lambda_n\). Its order is \(2^n\); the full linear
\(D_8\) coordinate is retained, even if its action on the vertex
quotient is not faithful. Its pair orders are \(4,4,2\). It has all
Boolean deletions and IP. Its third lower-central subgroup is
\(2\mathbb Z^2/\Lambda_n\), of order \(2^{n-5}\).

\emph{Proof.} The lattice indices are respectively \(2^{2e}\) and
\(2^{2e+1}\), so the normal form gives the stated orders. The three
linear pair products have orders four, four and two; hence the marked
pair orders cannot fall in either lattice quotient.

Write \(T=abcb\), \(U=cbab\), translations by \((-1,0)\), \((0,1)\). The
extra relators for \(H_n\) can be taken as \(T^{2^e},U^{2^e}\) in the
square case, and \(T^{2^e}U^{2^e},T^{2^e}U^{-2^e}\) in the diagonal
case. After deleting any marker, both \(T,U\) have order at most two:
deleting \(a\) leaves conjugates of \(c\), deleting \(c\) leaves
conjugates of \(a\), and deleting \(b\) leaves \(ac,ca\). Therefore
every lattice relator survives every deletion. The even Coxeter relators
also survive. Thus \(H_n\) has Boolean deletions and IP by UA.D.

Moreover, \[
 \gamma_3(W)=2\mathbb Z^2,
 \qquad \gamma_3(H_n)=2\mathbb Z^2/\Lambda_n,
 \qquad |\gamma_3(H_n)|=2^{n-5}.                         \tag{S2}
\] Indeed reduction modulo \(2\mathbb Z^2\) makes the adjacent squares
central; conversely \([[a,b],c]\) and \([[b,c],a]\) give the two doubled
unit translations, up to signs. Thus \(H_5\) is class two of order 32.
The squared translations are not always independent cyclic subgroups: in
\(H_6\) they coincide and have order two. Formula (S2) uses a lattice
index and does not require independence of those cyclic subgroups.
\(\square\)

\textbf{Lemma UA.SquareRecognition (all finite two-group square markings).}
Every finite two-group string C-group of type $\{4,4\}$ and order $2^n$
is marked-isomorphic to $H_n$; necessarily $n\ge5$.
\begin{proof}
Its two adjacent dihedral parabolics have orders eight and intersect in
order two, so their product already has 32 elements. Thus $n\ge5$.
The complete affine presentation of UA.Affine maps onto the given group $G$;
write $N$ for its kernel. IP gives irredundancy, and UA.Phi identifies
$G/G'$ with the elementary abelian group on its three markers. Summing
those three binary coordinates gives a character sending each reflection
to the nonidentity element of $C_2$. In the affine group this character is
orientation parity. Hence every element of $N$ preserves orientation.

The kernel contains no nonidentity torsion. A finite-order plane isometry
has a fixed point, obtained for example by averaging the orbit of any point.
The square reflection group moves that point into a closed fundamental right
triangle. The stabilizer of a point in that triangle is trivial in the
interior, generated by its wall reflection on an open edge, or the finite
dihedral group generated by the two walls meeting at a vertex. These
stabilizers are conjugates of the proper marked parabolics. Their maps into
$G$ are injective by its exact pair orders and IP, so none contains a
nonidentity element of $N$.

An orientation-preserving plane isometry is a translation or a rotation.
The linear part here has finite order in $D_8$, so a nontranslation rotation
has finite order. Thus all elements of $N$ are translations. Since $N$ is
normal of finite index, it is a $D_8$-invariant sublattice of $\Z^2$, of
index $2^{n-3}$. Under $\Z^2=\Z[i]$, invariance under a quarter-turn makes
it an ideal. The Euclidean algorithm in $\Z[i]$ makes it principal; its
generator has norm a power of two. As $2$ is a unit times $(1+i)^2$,
unique factorization shows that this generator is a unit times a power of
$1+i$. Even powers give the square lattice and odd powers the diagonal
lattice in UA.Square. There is one such ideal at each two-power index.
Therefore $N=\Lambda_n$ and the marked quotient is $H_n$.
\end{proof}

\textbf{Proposition UA.SquareDecomposition (the two square
decompositions).} For every integer \(q\ge1\), let \(M_2(q)\) and
\(M_1(q)\) be the groups obtained by adjoining respectively
\((bcba)^q=1\) and \((cba)^{2q}=1\) to the complete presentation (S1).
They are the square and diagonal lattice quotients, and have the
decompositions (S3) and (S4) below. Here \(D_{2q}\) is the dihedral
group of order \(2q\), including \(D_2=C_2\) and \(D_4=C_2^2\).

\emph{Proof.} Here bcba is translation by \((1,0)\), and \((cba)^2\) is
translation by \((1,1)\). Their signed-permutation orbits therefore
generate exactly the square and diagonal lattices in (S3)--(S4). This
identifies those complete presentations with the indicated lattice
quotients.

In the square case, \(\langle a,c^b\rangle\) acts on the first
coordinate and \(\langle a^b,c\rangle\) on the second. Their
translations have order q; each is \(D_{2q}\), since its reflection has
nonidentity linear part, also when q=1. They commute and intersect
trivially. The reflection b interchanges the two factors and lies
outside their product. They and b generate all three original
generators. Hence \[
 W/q\mathbb Z^2\cong(D_{2q}\times D_{2q})\rtimes C_2.
 \tag{S3}
\]

In the diagonal case, put \(A_0=\langle b^a,b^c\rangle\) and
\(B_0=\langle b,b^{ac}\rangle\). Their reflecting axes are respectively
\(x+y=0,1\) and \(x-y=0,-1\). The first factor has translations along
\((1,1)\), the second along \((1,-1)\), of order q in each case. Each
reflection retains its nonidentity linear part, so the two factors have
order 2q also at q=1. All cross-generators commute. Their linear parts
show that an element of their intersection must be a translation; the
independence of the two diagonal translations modulo the defining
lattice then makes that intersection trivial. Thus
\(A_0B_0\cong D_{2q}\times D_{2q}\).

Both a and c exchange the two factors by conjugation, so normalize their
product. The two individual reflections a,c do not belong to the product
by their linear parts. The only remaining possible nontrivial
intersection with \(\langle a,c\rangle\) is ac. Every translation part
of an element of \(A_0B_0\) with linear part \(-1\) has coordinates of
the same parity, whereas ac has translation part \((0,1)\). The diagonal
lattice preserves this parity condition, so ac is also outside the
product. Finally b already lies in \(B_0\), so the two factors and
\(\langle a,c\rangle\) generate the full group. Consequently \[
 W/q\langle(1,1),(1,-1)\rangle
 \cong(D_{2q}\times D_{2q})\rtimes(C_2\times C_2).
 \tag{S4}
\] These arguments specify the actions by the original affine generators
and prove the decompositions, rather than only equality of orders.
\(\square\)

\section{Changing the orders on selected pairs}\label{changing-the-orders-on-selected-pairs}

\textbf{Theorem UA.Edge (dihedral attachments).} Let
\(B=\langle b_i:i\in I\rangle\) be marked by nonidentity involutions and
have every Boolean deletion. Select a family \(\mathcal E\) of distinct
pairs \(S=\{i,j\}\) for which the product order \(q_S=|b_i b_j|\) is
finite and even. For each selected pair choose a positive multiple
\(m_S\) of \(q_S\), and attach the marked dihedral quotient \[
D_{2m_S}\longrightarrow B_S=D_{2q_S}.
\] The full compatible product of UA.SK is generated by its prescribed
marker tuples. It has the intersection property and all Boolean
deletions. Every selected pair has product order \(m_S\), and every
unselected pair retains its original product order. If \(B\) is finite,
then for every \(J\subseteq I\), \[
|H_J|=|B_J|\prod_{\substack{S\in\mathcal E\\S\subseteq J}}\frac{m_S}{q_S}.
\tag{E1}
\] If the base marking gives a regular hypertope, so does the
constructed marking, and its full type-preserving automorphism group is
\(H\).

\emph{Proof.} A group generated by two nonidentity involutions whose
product has finite order \(q_S\) is the dihedral group \(D_{2q_S}\); the
normal form alternates between powers of the product and a single
reflection. The indicated quotient exists because \(q_S\mid m_S\). It is
injective on each proper local parabolic, which is either trivial or
generated by one reflection. UA.D supplies IP in the base, and all input
deletions are part of the hypothesis. Thus UA.SK applies. The pair
supports form an antichain, so UA.SK.Local identifies every selected
pair subgroup with its full local dihedral group. UA.SK.Delete gives the
asserted deletions. The kernel on a selected pair has order \(m_S/q_S\),
and (E1) is UA.SK's order formula.

For an unselected pair, each extra coordinate of its product is a
reflection or the identity, because that coordinate's support cannot
contain both markers. If the base product has finite order \(q\), its
image under a deletion keeping just one of the two markers has order
two, so \(q\) is even. The extra coordinates therefore do not enlarge
its order. If the base product has infinite order, its lift has infinite
order by projection. This proves the assertion for every unselected
pair.

All local markers are dihedral reflections, hence involutions. Under the
final hypothesis, UA.SK.Geometry gives the geometric conclusion and
identifies the full type-preserving automorphism group. \(\square\)

\section{Amalgams and primary torsion}\label{amalgams-and-primary-torsion}

Fix injective homomorphisms \(i_A:C\to A\) and \(i_B:C\to B\). Let
\(A*B\) be the free product and let \[
P=(A*B)/\langle\!\langle i_A(c)i_B(c)^{-1}:c\in C\rangle\!\rangle.
\tag{A1}
\] We use the amalgam normal-form theorem in the following precise form:
the canonical maps of \(A\) and \(B\) into \(P\) are injective; an
alternating reduced word with factors outside the common copy of \(C\)
has nonzero reduced length and is nonidentity; every element is
conjugate either into a factor or to a cyclically reduced alternating
word of length at least two, whose positive powers remain reduced with
proportionally increasing length. This is a retained classical theorem.

\textbf{Theorem UA.Amalgam (universal property and torsion).} The group
\(P\) in (A1), with its canonical maps \(j_A,j_B\), is a pushout of
\(A\leftarrow C\to B\) in the category of groups: for every group \(D\)
and homomorphisms \(f_A:A\to D\), \(f_B:B\to D\) satisfying
\(f_Ai_A=f_Bi_B\), there is a unique homomorphism \(f:P\to D\) with
\(fj_A=f_A\) and \(fj_B=f_B\).

Let \(p\) be a prime. Say that a group has \(p\)-primary finite torsion
when every finite-order element has \(p\)-power order. This property
passes to subgroups, to extensions with finite \(p\)-group quotient, and
to amalgamated free products of groups having the property. If
\(i_A(C)\ne A\) and \(i_B(C)\ne B\), then \(P\) is infinite. If \(A\)
and \(B\) are finitely generated, then so is \(P\).

\emph{Proof.} The free-product universal property gives a unique
homomorphism \(A*B\to D\) extending \(f_A,f_B\). Agreement on \(C\) puts
every defining relator in its kernel, so it factors through \(P\). The
images of \(A\) and \(B\) generate \(P\); hence this factorization is
unique. The equality \(j_Ai_A=j_Bi_B\) follows from the defining
relations, completing the pushout property.

The torsion assertion for subgroups is immediate. For an extension
\(1\to K\to E\to R\to1\), with \(R\) a finite \(p\)-group and \(K\)
having \(p\)-primary finite torsion, take any finite cyclic subgroup
\(\langle x\rangle\) of \(E\). Its intersection with \(K\) is a cyclic
\(p\)-group, and its image in \(R\) is a \(p\)-group. The order formula
for this restricted exact sequence makes \(|x|\) a power of \(p\).

For an amalgam, the normal-form theorem places a torsion element, up to
conjugacy, in a factor: the alternative cyclically reduced word has
powers of unbounded reduced length and cannot have finite order. The
factor hypothesis therefore proves the torsion assertion. If the common
subgroup is proper in both factors, choose \(a\in A\setminus i_A(C)\)
and \(b\in B\setminus i_B(C)\). The words \((ab)^n\) are reduced of
length \(2n\), so are all distinct; this proves infinitude. Finally, the
images of finite generating sets of \(A\) and \(B\) generate their free
product and its quotient \(P\), proving finite generation. \(\square\)

The simultaneous fiber product of UA.SK has its displayed compatibility
equations and its proved generation theorem. It is not the pushout in
UA.Amalgam. A graph of proof dependencies is neither of these group
constructions.

\section{Recognizing an affine presentation}\label{recognizing-an-affine-presentation}

We write an integral affine transformation as \((M,t)\), acting by
\(v\mapsto Mv+t\), with multiplication

\[ (M,t)(N,u)=(MN,t+Mu). \tag{1} \]

Let \(E\) be a group generated by finitely many such transformations, and
let \(Q\leq\operatorname{GL}_d(\mathbb Z)\) be its finite linear image.
Its translation kernel \(L\) is a subgroup of \(\mathbb Z^d\), invariant under
\(Q\). A translation is written \(t_v=(I,v)\).

\textbf{Lemma AL.1 (fixed Schreier lattice).} Choose a word lift \(\widehat q\in E\)
for each \(q\in Q\), with \(\widehat1=1\). Then \(L\) is generated by the
translations

\[\widehat{\overline s q}^{-1}s\widehat q,\qquad
s\text{ a distinguished generator},\quad q\in Q.\tag{2}\]

\emph{Proof.} The products have trivial linear image. Write a partial word
as \(\widehat q n\), with \(n\in L\). Left multiplication by \(s\) changes
this to \(\widehat{\overline s q}t n\), where \(t\) is the transition in
(2). Induction expresses a word of trivial linear image as a product
of these transitions; its final representative is \(\widehat1=1\).
Inverses introduce no extra generator: the transition for $s^{-1}$ at $q$ is the inverse of the transition for $s$ at $\overline{s}^{-1}q$. Thus the induction applies to words with positive and negative letters. Every transition itself is in \(L\), so they generate exactly \(L\).
\(\square\)

The normalization in (2) matters. The raw differences
\(s\widehat q\widehat{\overline s q}^{-1}\) need not generate \(L\)
without conjugation. Each is a conjugate of the corresponding transition
by \(\widehat{\overline s q}\). Consequently their \(Q\)-orbits generate
\(L\). In particular, if a proposed \(Q\)-invariant lattice contains every
raw difference, it contains every normalized transition and hence \(L\).
If exhibited raw differences also span the proposed lattice, equality
follows. This invariant-lattice version justifies the raw-difference
format of the unchanged earlier \(\{6,6\}\) checker {[}C66{]}. The new
\(\{6,4\}\) checker uses (2) directly.

This is a finite proof of a lattice identity. Containment in a proposed
lattice is checked on the normalized transitions, or on the raw differences
with the stated invariance condition. Equality follows when independent
integer combinations give a basis of that lattice, or when the greatest
common divisor of the full-rank minors gives its exact index. Sampling
finite quotients is not a replacement for this step.

\textbf{Lemma AL.2 (presentation bridge).} Let \(U\) be a finitely presented marked
group whose relators hold on the affine generators of \(E\). Suppose their
linear images define the same epimorphism to \(Q\). If an exact
Reidemeister--Schreier presentation proves
\(\ker(U\to Q)\cong\mathbb Z^d\), and the affine translation kernel \(L\)
has rank \(d\), then the marked epimorphism \(U\to E\) is an isomorphism.

\emph{Proof.} The restriction maps \(\ker(U\to Q)\) onto \(L\): an affine
translation has a preimage in \(U\), and its linear image is trivial.
This is a surjection between free abelian groups of rank \(d\). Tensoring
with \(\mathbb Q\) shows that its kernel has rank zero. A subgroup of a
free abelian group is torsion-free, so that kernel is trivial. The
remaining kernel of \(U\to E\) lies in this subgroup, and is therefore
trivial as well. \(\square\)

The hypothesis concerns the actual kernel presentation. Computing only
its abelianization would not establish it: a nonabelian group can have
free abelian abelianization of rank \(d\). Our applications retain the
exact rewriting and Tietze algorithms used to verify this finite
presentation step. It is a computer-assisted input, not a Lean proof.

\section{All parabolic intersections in an affine period quotient}
\label{sec:affine-parabolics}

The finite adjacent subgroups in the polyhedral applications have faithful
linear images. An arbitrary parabolic need not have that property: it can
contain translations. We therefore state the intersection calculation before
making the faithful-image specialization. It uses the actual affine group,
whose identification with a presented group is the separate obligation in
AL.2.

Let $E=\langle r_i:i\in I\rangle\leq\mathbb Z^d\rtimes Q$, where $I$ and
the linear image $Q$ are finite. Multiplication is
$(q,u)(q',v)=(qq',u+qv)$. Write $L\leq\mathbb Z^d$ for the translation
kernel of $E$. For $J\subseteq I$ put
\[
 E_J=\langle r_i:i\in J\rangle,\qquad
 Q_J=\operatorname{im}(E_J\to Q),\qquad
 K_J=\{v:(1,v)\in E_J\}\leq L.
\]
Choose, for each $q\in Q_J$, an element $(q,u_J(q))\in E_J$, with
$u_J(1)=0$. These are finite choices of affine word lifts, not assumptions
about a splitting of $E_J$. The entire fiber over $q$ is
\begin{equation}
 \{(q,u):u\in u_J(q)+K_J\}.\label{eq:affine-fiber}
\end{equation}
Indeed the quotient of two elements with the same linear part is their
translation difference; conversely left multiplication by $(1,k)$ adds any
$k\in K_J$ to the translation part.

Let $M\leq L$ be $Q$-invariant. It is then a normal translation subgroup of
$E$. Write $\overline E=E/M$ and $\overline E_J$ for the image of $E_J$.
The subgroup $M$ can have any rank, including zero.

\begin{theorem}[AA.ParabolicPeriods]\label{thm:all-periods}
The marking of $E/M$ has the intersection property if and only if, for
every $J,T\subseteq I$, with $C=J\cap T$, both conditions hold:
\begin{align}
 (K_J+M)\cap(K_T+M)&=K_C+M,\label{eq:period-kernels}\\
 u_J(q)-u_T(q)&\notin K_J+K_T+M
 \quad(q\in(Q_J\cap Q_T)\setminus Q_C).\label{eq:period-exclusions}
\end{align}
More generally, a common linear part $q\in Q_J\cap Q_T$ occurs in
$\overline E_J\cap\overline E_T$ exactly when
\[
 u_J(q)-u_T(q)\in K_J+K_T+M.
\]
When it occurs, its translation fiber is a coset of
$((K_J+M)\cap(K_T+M))/M$.
\end{theorem}
\begin{proof}
Modulo $M$, equation~\eqref{eq:affine-fiber} becomes the coset
$(u_J(q)+K_J+M)/M$. Equality in $E/M$ preserves the linear part, since
$M$ consists of translations. Two cosets of additive subgroups intersect
if and only if the difference of their representatives belongs to the sum
of those subgroups. When nonempty, their intersection is a coset of the
intersection of the subgroups. This proves the two fiber assertions.

For $q\in Q_C$, a lift $(q,u_C(q))\in E_C$ belongs to both $E_J$ and
$E_T$. It can therefore be used as a common representative of their two
translation cosets. The intersection fiber is
\[
 \bigl(u_C(q)+(K_J+M)\cap(K_T+M)\bigr)/M,
\]
whereas the fiber of $\overline E_C$ is
$(u_C(q)+K_C+M)/M$. Condition~\eqref{eq:period-kernels} makes them equal.
Condition~\eqref{eq:period-exclusions} excludes every other possible linear
part. Together these prove
$\overline E_J\cap\overline E_T=\overline E_C$.

Conversely that equality, at the identity linear part, implies
\eqref{eq:period-kernels}: the relevant additive groups all contain $M$, so
equality of their quotients by $M$ is equality of the groups themselves.
For a linear part outside $Q_C$ its intersection fiber must be empty,
which gives~\eqref{eq:period-exclusions}. Applying this to every pair of
index sets proves the equivalence. Changing any chosen lift adds an element
of $K_J$ or $K_T$ and does not change the membership test.
\end{proof}

All data in this criterion can be obtained by finite exact calculations
for a fixed period lattice. Enumerate $Q_J$ from the linear generator
matrices, retain a word lift of each element, and apply AL.1 to $E_J$ and
$Q_J$. Its normalized Schreier translations generate exactly $K_J$.
Integer lattice sum, intersection and membership calculations then decide
\eqref{eq:period-kernels}--\eqref{eq:period-exclusions}. No finiteness of
$E_J$ and no injectivity of $E_J\to Q_J$ is assumed. For a parameterized
family $M_m$ the same equations remain exact, but an all-parameter solution
of them still needs a proof; finitely many successful periods do not
establish it.

The first condition cannot be omitted when parabolics contain translations.
Take $E=\mathbb Z^2$ with markers $e_1,e_2$ and trivial $Q$. For
$M=\langle e_1-e_2,2e_2\rangle$, the two distinct rank-one parabolics have
the same nontrivial image in $E/M\cong C_2$, so IP fails. Here
$K_{\{1\}}+M=K_{\{2\}}+M=\mathbb Z^2$, whereas $K_\varnothing+M=M$.
There is no nonidentity linear part to detect this failure. By contrast,
$M=2\mathbb Z^2$ gives the two distinct coordinate factors of $C_2^2$ and
satisfies the criterion. These examples concern general marked groups;
the theorem does not require involutory markers.

This statement extends the previously proved rank-three period argument to
all marked parabolics. Its proof is the classical intersection formula for
cosets of additive subgroups. The extension and the explicit hypothesis
separation are part of this edition; no claim of historical priority is
made. Existing formal period certificates cover the previously stated
faithful-parabolic monotonicity, not this full arbitrary-parabolic theorem.

\section{Periods and intersections of distinguished subgroups}\label{periods-and-intersections-of-distinguished-subgroups}

Let \(w_1,\ldots,w_d\in U=E\) be translation words with linearly independent
vectors \(v_1,\ldots,v_d\). Write
\(\Lambda=\sum_i\mathbb Z v_i\leq L\) and suppose it is \(Q\)-invariant.
For a positive integer \(m\), adjoining \(w_i^m=1\) to the complete
presentation of \(U\) gives exactly

\[ E_m=E/(m\Lambda).\tag{3}\]

Indeed, the normal closure of the powers is \(m\Lambda\): conjugation by
an affine element acts on translations by its linear part, and this
lattice is already invariant. Hence

\[ |E_m|=|Q|[L:\Lambda]m^d.\tag{4}\]

This proves equality for the presented group, not merely the order of a
chosen quotient. If \(Q\) is solvable, every \(E_m\) is solvable.

Suppose now that the marking has rank three with distinct nonidentity involutions,
its endpoints commute, both adjacent products have order at least three, and
its two adjacent subgroups \(A=E_{01}\) and \(B=E_{12}\) are finite and
inject into \(Q\). For each \(q\in\overline A\cap\overline B\), let
\((q,t_A(q))\) and \((q,t_B(q))\) be their unique affine lifts. Put

\[ \Delta(q)=t_A(q)-t_B(q).\tag{5} \]

\textbf{Theorem AL.3 (period-lattice intersection criterion).} The images of \(A\)
and \(B\) in \(E_m\) meet precisely in the elements indexed by

\[ \{q\in\overline A\cap\overline B:\Delta(q)\in m\Lambda\}.\tag{6}\]

In particular, if \(A\cap B=\langle r_1\rangle\) and every other value
\(\Delta(q)\) lies outside \(\Lambda\), then every \(E_m\), \(m\geq1\), has
the full rank-three intersection property and the same adjacent orders
as \(E\). Its order is (4).

\emph{Proof.} Apply the fiber assertion of AA.ParabolicPeriods to
$J=\{0,1\}$, $T=\{1,2\}$ and $M=m\Lambda$. Faithful linear images give
$K_J=K_T=0$, so its general coset-membership condition is exactly (6).
They also show that both adjacent subgroups remain injective in the
quotient. Since $m\Lambda\subseteq\Lambda$, the stated exclusions leave
only the common involution. The rank-three test gives all subset
intersections. Adjacent orders and (4) follow from the injectivity and
the full translation index already established. \(\square\)

The proof separates two possibilities often hidden by equal group-order
counts. The point-group intersection may already have order two, as in
the \(\{6,4\}\) application. Or it may be larger, with a short list of
nonzero translation differences, as in the \(\{6,6\}\) application. Those
differences must be excluded from the period lattice explicitly.

\subsection*{One coarse quotient controls all smaller periods}
The same argument has a group-theoretic formulation which is useful independently of coordinates. Let $A,B,L$ be subgroups of a group $E$, with
\[
 L\cap A=L\cap B=1.
\]
Let $N\unlhd E$ lie in $L$. If the images of $A,B$ in $E/N$ intersect exactly in the image of $A\cap B$, then the same statement holds in $E/M$ for every normal subgroup $M\le N$.

Indeed, suppose $a\in A$, $b\in B$ have the same image modulo $M$. They have the same image modulo $N$, so there is $c\in A\cap B$ with $aN=cN=bN$. Then $ac^{-1}\in N\cap A=1$ and $bc^{-1}\in N\cap B=1$, proving $a=b=c$. The reverse inclusion is immediate. In the affine setting $L$ is the translation kernel, $N=\Lambda$ and $M=m\Lambda$. Point-faithfulness supplies precisely the two trivial intersections. This statement, for an arbitrary indexed family of smaller normal periods, is verified in \texttt{AffinePeriodIntersections.lean}. The finite coordinate checks which identify $L,\Lambda$ and the initial overlap remain separate mathematical obligations.

\section{Changing one adjacent order}\label{changing-one-adjacent-order}

Let \((B;a,b,c)\) be a rank-three string C-group of type \(\{h,q\}\), with
\(h,q\) even, and suppose the map

\[\rho:B\longrightarrow B_{12},\qquad a\mapsto1, b\mapsto b, c\mapsto c
\tag{7}\]

is a retraction. For \(k\) a positive multiple of \(q\), take the standard
marked epimorphism \(D_{2k}\to B_{12}=D_{2q}\). The compatible product

\[F=\{(g,d)\in B\times D_{2k}:\rho(g)=\overline d\}\tag{8}\]

has markers \((a,1),(b,\widetilde b),(c,\widetilde c)\). The evenness of
\(h\) is forced by (7); it also ensures that the first pair has unchanged
order \(h\).

\textbf{Theorem AL.4 (one-edge extension).} These markers generate \(F\). If \(B\) is finite, its order
is \(|B|k/q\). Its type is \(\{h,k\}\), and it has the intersection
property. The projection onto \(B\) is injective on the first adjacent
subgroup. If \(B\) is solvable, so is \(F\).

\emph{Proof.} Apply UA.SK to the single support $S=\{1,2\}$, using the given deletion and the marked dihedral epimorphism $D_{2k}\to D_{2q}$. Its restrictions to both reflection subgroups are injective, which verifies every proper local subgroup. Its kernel has order $k/q$. The theorem gives generation of the whole compatible product, its order and IP. UA.SK.Local identifies the second adjacent subgroup with $D_{2k}$; the first has trivial projection kernel by (K2), since it does not contain $S$. Thus its order and adjacent product order remain unchanged. Explicitly the other local coordinate of that product has order two, which divides the even integer $h$. The cyclic projection kernel and the solvability of the base prove solvability. \(\square\)

Theorem AL.4 is a specialization of UA.SK; it is included to make clear that
only the indicated deletion is needed here. More simultaneous edges use
the supported coordinate kernels of that theorem.

There is a useful exact presentation variant. Suppose a complete marked
presentation for \(B\) contains \((bc)^q=1\). Replace that relation by
\((bc)^k=1\), retain the remaining relators, and add
\([a,(bc)^q]=1\). Call the resulting presented group \(H\). If every retained
relator holds in (8), then \(H=F\) as a marked group. Here the local dihedral presentation holds in $H$ by the involution relations and $(bc)^k=1$. Its kernel word $(bc)^q$ generates a normal subgroup: $b,c$ invert it and the new relation makes $a$ centralize it. Killing that word gives exactly the original marked presentation of $B$. These are all four conditions of UA.Recognition, so $H=F$ even when the base is infinite. No finite order comparison is needed.

\section{Recognizing supported extensions of affine quotients}
\label{sec:period-extensions}

The following interface combines the two recognition steps. It separates a
presentation of the affine base from the presentation of a group with larger
local sections. The latter is not identified merely by displaying a finite
image of the right type.

\begin{theorem}[AA.PeriodExtension]\label{thm:period-extension}
Let $I$ be finite and let $U=\langle r_i:i\in I\rangle$ be a finitely
presented marked group with an affine marked image
$E\leq\mathbb Z^d\rtimes Q$, where the finite linear images agree. Suppose
the actual kernel of $U\to Q$ has a presentation identifying it with
$\mathbb Z^d$, and that the translation kernel $L$ of $E$ has rank $d$.
Let $M\leq L$ be $Q$-invariant, with an explicitly given finite lattice
basis represented by words in $U$.

Suppose the equations of AA.ParabolicPeriods hold for $M$. Put $B=E/M$.
Choose a finite antichain $\mathcal F$ of label supports. For each
$S\in\mathcal F$, suppose there is a deletion $p_S:B\to B_S$ and a
marked epimorphism $q_S:V_S\to B_S$ from a group generated by the
markers on $S$, injective on every proper marked
subgroup of $V_S$. Write $N_S=\ker q_S$.

Let $P$ be a marked presentation with the following properties:
\begin{enumerate}
\item Its relators hold in the compatible product of $B$ and the $V_S$
with respect to $p_S,q_S$.
\item Every local presentation for $V_S$ holds in $P$.
\item The image in $P$ of each $N_S$ is normal in $P$.
\item Killing these local kernel images leaves exactly the presentation
of $U$ with the basis words of $M$ set equal to one.
\end{enumerate}
Then $P$ is marked-isomorphic to that full compatible product. It has IP,
and its parabolic projection kernels are the coordinate products
\[
 \ker(P_J\to B_J)=\prod_{S\subseteq J}N_S.
\]
Each local subgroup $P_S$ is exactly $V_S$. If $[L:M]$ and all $N_S$ are
finite, then
\[
 |P|=|Q|[L:M]\prod_{S\in\mathcal F}|N_S|,
 \qquad |P_J|=|B_J|\prod_{S\subseteq J}|N_S|.
\]
The geometric and deletion conclusions hold precisely under the additional
hypotheses stated in UA.SK.Geometry and UA.SK.Delete, respectively.
\end{theorem}
\begin{proof}
AL.2 identifies $U$ with $E$ using the actual free abelian kernel
presentation and its surjection onto the rank-$d$ translation lattice.
The normal closure of the chosen basis translations of $M$ in $E$ is
$M$: it contains their additive span, and invariance under the linear
image shows that this span is already normal. Thus adjoining these words
identifies the whole base presentation with $B$, not only with a quotient
of $B$. AA.ParabolicPeriods gives IP in $B$.

The stipulated deletions and proper-parabolic injectivity are exactly the
hypotheses of UA.SK. That theorem constructs and generates the full
compatible product, identifies every projection kernel, and proves IP.
The four presentation hypotheses are exactly UA.Recognition. Its proof uses
normality and the faithful separate local coordinates to exclude every
remaining kernel; it proves $P$ is that compatible product. The antichain
clause UA.SK.Local identifies its local subgroups. When the displayed
indices are finite, $|B|=|Q|[L:M]$, and UA.SK gives the stated orders.
The last sentence invokes the additional proved hypotheses of the two
named corollaries; IP alone does not establish a hypertope's flag property.
\end{proof}

The theorem also applies with no local supports, in which case $P=B$.
For an affine base with an elementary normal-form proof, that proof can
replace the Schreier computation in the first paragraph: the required input
is the exact marked isomorphism $U=E$, not a particular algorithm for proving
it. This is how the square groups enter the two-power applications. The
cubic $\{6,4\}$ and $\{6,6\}$ bases instead retain their two full exact
kernel-presentation programs. Their lattice and overlap programs verify the
separate arithmetic hypotheses. No one of these certificates replaces the
others.

This theorem assembles already proved normal-kernel and supported-extension
arguments with the period criterion. The eight applications below retain
their source authorship. The high-type and order-1024 classifications also
need the forest, cyclic-core and exhaustive-kernel proofs; they are not
consequences of this construction theorem alone.

\section{Retained inputs, applications and attribution}\label{retained-inputs-applications-and-attribution}

The proofs use the elementary subgroup, quotient and isomorphism
theorems. Square recognition also uses the Euclidean algorithm and unique
factorization in $\Z[i]$, and the elementary point stabilizers in the
square reflection tiling described in its proof. The further classical
inputs are explicitly limited:
Burnside's basis theorem for finite \(p\)-groups; their nilpotence;
generation of the lower-central graded Lie ring by degree one and its
multilinear commutator rules; and the precise amalgam normal-form
theorem stated in the amalgam section. The Frattini, forest, finite-quotient and
torsion deductions are proved in full above. No finite group
classification, atlas census or numerical search result supplies a
premise to these theorems.

For applications to abstract regular polytopes we retain the following
standard correspondence {[}MS02, Section 2E{]}. A group generated by an
ordered list of nonidentity involutions, with the string commutation
rules and IP, is the full automorphism group of a regular abstract
polytope with that distinguished marking. Its rank is the length of the
list, its type consists of the adjacent product orders, and its number
of flags is the group order. A generator-preserving epimorphism between
such groups induces a rank- and adjacency-preserving covering of their
associated polytopes. These are geometric input theorems, not new
conclusions of the present paper. By contrast, the coset-geometric
recognition and coordinate-kernel Helly statements needed for the
hypertope branch are proved in the incidence section with their full hypotheses.

The supported-extension construction supplies a common proof interface
for the two-power constructions studied in {[}HFL20, HFL19a, HFL19b,
CFHS25{]}. The presentation-recognition theorem supplies the additional normality and
exact-base conditions needed to identify a compatible product with a complete
universal presentation; local construction alone does not do this.
The forest quotient gives the canonical group underlying the main
quotient and cover conclusions of {[}HS26{]}, and the finite-kernel
parameter space gives the exhaustive part of the order-1024
reconstruction associated with {[}GLP18{]}. Their exact source
statements, literal corrections, numerical certificates and named-object
identifications belong to the separate applications.

The supported-kernel induction, forest quotient and affine models are
reused complete proofs. The normal-kernel recognition theorem replaces
repeated finite-order comparisons, and the mix criterion separates joint
coordinate detection from generation of a fiber product. Every application
still proves its local defining relations, normality and base identifications.
Original source results and borrowed geometric theorems keep their attribution.

Existing Lean certificates cover declared components of the kernel,
supported-extension and Helly arguments. They do not presently formalize
this whole paper: in particular, the full forest quotient proof, the
affine model and the geometric recognition/application interfaces are
not all connected in Lean. Their written proofs above remain part of the
mathematical obligation. The accompanying certificate scope
distinguishes inherited checked declarations from any freshly checked
application components.


The mixed-prime application also retains Burnside's $p^a q^b$ solvability
theorem, the exact Lucchini cyclic-core bound as quoted with its hypotheses,
and the complete SmallGroups catalogue at order 384. The affine bridge
uses actual Reidemeister--Schreier and Tietze computations of two kernels,
not their abelianizations. The graph-symmetry application retains Frucht's
theorem and the complete SmallGroups catalogues at orders 40,48,56 for its
isolated rank-three minimum. These are application inputs; none is silently
used to prove the general lattice-coset or normal-kernel theorems.

All eleven complete required program/data files are retained once in this
volume: four files for the two-group classification (including its complete
named-presentation data and runner), six for the mixed-prime affine families,
and one for the graph-symmetry minimum. The three complete appendices print
those bytes. Operational drivers and all previously scoped Lean files are
separately inventoried. Existing successful semantic and computational
reviews are reused for unchanged bytes; this consolidation does not claim
to have reexecuted every old census or to have formally verified the new
arbitrary-parabolic lattice criterion.

\section*{References}
{[}MS02{]} P. McMullen and E. Schulte, \emph{Abstract Regular
Polytopes}, Cambridge University Press, 2002. The regular-polytope and
covering correspondences are retained at the scope stated above.

{[}HFL20{]} D.-D. Hou, Y.-Q. Feng and D. Leemans, \emph{On regular
polytopes of 2-power order}, revised author manuscript dated 8 November
2019; published in \emph{Discrete \& Computational Geometry} 64 (2020),
339--346. DOI: 10.1007/s00454-019-00119-5.

{[}HFL19a{]} D.-D. Hou, Y.-Q. Feng and D. Leemans, \emph{Existence of
regular 3-polytopes of order \(2^n\)}, arXiv:1802.08947v2, 22 January
2019; published in \emph{Journal of Group Theory} 22 (2019), 579--616.

{[}HFL19b{]} D.-D. Hou, Y.-Q. Feng and D. Leemans, \emph{Existence of
regular 3-hypertopes with \(2^n\) chambers}, arXiv:1901.09034v1, 23
January 2019; published in \emph{Discrete Mathematics} 342 (2019),
1857--1863.

{[}HS26{]} D.-D. Hou and E. Schulte, \emph{String C-groups of 2-power
order project onto a common string C-group}, arXiv:2605.20691v1, PDF
dated 21 May 2026. The specified source edition is used without
asserting that every auxiliary claim is correct.

{[}CFHS25{]} G. Cunningham, Y.-Q. Feng, D.-D. Hou and E. Schulte,
\emph{Symmetric polytopes whose automorphism groups are 2-groups},
arXiv:2512.15511v1, PDF dated 18 December 2025.

{[}GLP18{]} Y. Gomi, M. L. Loyola and M. L. A. N. De Las Peñas,
\emph{String C-groups of order 1024}, \emph{Contributions to Discrete
Mathematics} 13(1) (2018), 1--22. DOI: 10.55016/ojs/cdm.v13i1.62475.

{[}HFL24{]} D.-D. Hou, Y.-Q. Feng and D. Leemans,
\emph{Regular 3-polytopes of order $2^n p$},\hfill\break arXiv:2001.02945v2,
14 July 2024, 17 pages. The application preserves that exact edition's
period-sign correction and all four principal results.

{[}LST26{]} D. Leemans, P. Spiga and P. Tranchida,
\emph{Hypertopes with prescribed diagram symmetries},
arXiv:2610.09711v1 (2026). Its graph construction is the seventh
previously admitted Topic A application.
\end{document}
